\section{The added Jacobian, absolute-curve, and prime-orbit papers}
\label{sec:sites-followup}

\subsection{Exact source inventory and the dependence on A}

The following inventory concerns the supplied source versions, not a
claim about a later corrected version.  Source labels below are the
literal LaTeX labels; where an environment has no label, its section and
opening words identify it.  An assertion appearing as a theorem in a
source is distinguished from an assertion verified below.
\begin{itemize}
\item Connes--Consani, \emph{The Jacobian of $\operatorname{Spec}\mathbb Z$}
(2602.15941), \texttt{Jacobian.tex}: Theorem \texttt{xqthm1} in
\S\texttt{adrel} identifies the Picard monoid of rank-one arithmetic
divisors with
$X_{\mathbb Q}=\mathbb Q^\times\backslash\mathbb A_{\mathbb Q}/
\widehat{\mathbb Z}^{\times}$.  Section \texttt{jacob} defines
$\operatorname{Jac}(\overline{\operatorname{Spec}\mathbb Z})
=\operatorname{Pic}(\overline{\operatorname{Spec}\mathbb Z})/
\mathbb R_+^\times$; equation \texttt{extAJ} specifies the Abel--Jacobi
map.  Theorem \texttt{unif} identifies the adelic monoid
$Y_{\mathbb Q}=\mathbb Q^\times\backslash\mathbb A_{\mathbb Q}$ with
the monoids of framed and rooted arithmetic divisors.  These are
unconditional constructions.  The finite-prime fiber theorem in
\S\texttt{fibers} gives $C_{\mathbb Q}/\mathbb Q_p^\times$.
The archimedean entry of the introductory Theorem \texttt{thm:fibers}
and the unlabelled theorem in ``The fiber over the archimedean orbit''
require the sign correction proved in AC-G13 below.
\item The same paper, \S\texttt{Wtrace}, equation
\texttt{weilexplicit}, restates the Weil explicit formula for all
compact idele-class characters.  Its \S\texttt{geomtrace} explains the
factor $|1-u|_v^{-1}$ by the transverse scaling kernel; the computation
at $u=1$ still needs the regularization used in the explicit formula.
Equation \texttt{semilocTrace} is the already established semilocal
cutoff trace formula, including its divergent term
$2h(1)\log\lambda$.  The final unlabelled theorem in
\S\texttt{sectshespec} recalls a sheaf of semilocal crossed products.
None of these constructions assumes A.  A locates their Dirichlet
spectral zeros and gives the Hermitian positivity discussed in AC-G01;
it does not remove the cutoff or provide a new surface cohomology.
\item Connes--Consani, \emph{Absolute geometry of
$\operatorname{Spec}\mathbb Z$} (2606.06604), \texttt{FF.tex}:
Proposition \texttt{prop12} identifies the pulled-back stalk at $p$ as
$\mathbb F_1[T^{\mathbb Z[1/p]_+}]$.  Theorem \texttt{main} identifies
its local points in a perfectoid field $C$ of residue characteristic $p$
with $1+\mathfrak m_{C^\flat}$, equivariantly for
$\mathbb Q_p^\times$.  Proposition \texttt{lurieprop2} and Corollaries
\texttt{luriecor3}, \texttt{luriecor4}, in the explicitly
algebraically-closed characteristic-$p$ setting of
\S\texttt{ffcurve}, identify the quotients with untilts and the closed
points of the Fargues--Fontaine curve.  Theorem
\texttt{scholze\_heuristic} and the unlabelled theorem in ``The universal
tilting functor and the geometric sieve'' give the corresponding
matched-prime and mismatched-prime descriptions.  Algebraic closedness,
the topology, and the distinction between trivial and nontrivial points
are part of these statements, not consequences of A.
\item In \texttt{FF.tex}, Propositions
\path{prop:archimedean_points_C} and
\path{prop:nonarchimedean_points_C} describe the complex-valued local
points.  Theorem \texttt{thm:tate\_decomposition} in ``Rectangular
decomposition of $E_p$'' gives
$E_p=\mathbb C^\times/p^{\mathbb Z}
\simeq C_p\times\widetilde{\mathcal X}_\infty$ as a product of
\emph{real} one-dimensional manifolds, with the complex structure
specified by $d\lambda/\lambda+i\,d\theta$.  This source label also
occurs in the introduction, so the section is necessary to disambiguate
it.  The period ratio is $i\log p/(2\pi)$ and
$E_p(\mathbb R)=\mathbb R^\times/p^{\mathbb Z}$ has two components.
These statements do not assume RH, do not construct a common elliptic
curve over $\mathbb Q$ with these reductions, and do not identify a
global arithmetic-square intersection theory.  The final paragraph of
\S\texttt{outlook} explicitly leaves the descent of Frobenius
eigenspaces such as $B^{\varphi=p}$ and their relation to $C_p$ for
future work.
\item Connes--Consani, \emph{Knots, primes and the adele class space}
(2401.08401), \texttt{sect2.tex}, equation \texttt{inverseim}, proves
the mapping-torus description
\[
 \pi^{-1}(C_p)=
 \left(\mathbb R_+^\times\times
        \prod_{q\ne p}\mathbb Z_q^\times\right)/p^{\mathbb Z}.
\]
Theorems \texttt{mainglobalintro}, \texttt{mainintro} identify its
monodromy with the image of arithmetic Frobenius in the appropriate
\emph{abelianized} etale fundamental group.  Proposition
\texttt{ergodic23} of \texttt{sect3.tex} proves the three-place
stratification.  Proposition \texttt{bciso} of \texttt{sect4.tex}
constructs a locally compact three-dimensional classifying space with
a codimension-one lamination; the source expressly says this dimension
comes from choosing two finite primes, not from the homological
dimension of $\operatorname{Spec}\mathbb Z$.  All are unconditional.
\item Connes--Consani, \emph{Knots, primes and class field theory}
(2501.06560), \texttt{arxiv-pisa.tex}: Definition \texttt{coverdef}
defines $X_{\mathbb Q}^{\chi}$ for a continuous surjection
$\chi:\widehat{\mathbb Z}^\times\to G$ to a finite abelian group.
Lemma \texttt{Gacts} identifies the inertia stabilizer; Theorem
\texttt{mappingtorus1} identifies unramified monodromy with
$\chi(p)$; Theorem \texttt{main} constructs the class-field functor;
Theorem \texttt{main2} identifies the maximal abelianized
etale monodromy.  Theorem \texttt{structure1} gives the crossed-product
sheaf, whose generic stalk is
$\mathcal S(\mathbb A_{\mathbb Q})\rtimes\mathbb Q^\times$ and whose
global sections are $\mathcal S(\mathbb R)\rtimes\{\pm1\}$.
These results are unconditional.  The trivial-group exception in
Corollaries \texttt{ramicor}, \texttt{monocor} and Theorem
\texttt{main}(ii) is addressed in AC-G17.
\end{itemize}

\subsection{The sign at infinity and the actual Jacobian monoid}

\begin{proposition}[AC-G13: corrected archimedean fiber]
\label{AC-G13}
Use exactly the quotient
$Y_{\mathbb Q}=\mathbb Q^\times\backslash\mathbb A_{\mathbb Q}$
in \texttt{Jacobian.tex}, and put $e_\infty=(1_f,0_\infty)$.
The fiber above the Abel--Jacobi point
$\Theta(\infty)=\mathbb Z_\infty$ is canonically
\[
 \mathcal F_\infty
 \simeq C_{\mathbb Q}/\iota_\infty(\mathbb R^\times)
 \simeq\widehat{\mathbb Z}^{\times}/\{\pm1\}.
\]
Under the cyclotomic Artin identification its Galois interpretation is
$\operatorname{Gal}(\mathbb Q^{\mathrm{ab}}\cap\mathbb R/\mathbb Q)$.
The quotient by $\mathbb R_+^\times$ asserted in the supplied source is
too small.  In particular the sign character, and every odd Dirichlet
character, does not descend to a character of this archimedean fiber.
\end{proposition}
\begin{proof}
The fiber consists of classes of adeles whose archimedean component is
zero and whose associated rank-one group is isomorphic to $\mathbb Z$.
Rescaling by a rational number puts the finite components in
$\widehat{\mathbb Z}^{\times}$.  Two such normalized representatives
$a_f,b_f$ define the same class precisely when $b_f=q a_f$ for
$q\in\mathbb Q^\times$ with $v_p(q)=0$ at every finite prime.
Such a rational number is $1$ or $-1$.  Hence the fiber is
$\widehat{\mathbb Z}^{\times}/\{\pm1\}$, with the quotient topology
and componentwise group law.

Alternatively the map $u\mapsto[u e_\infty]$ from the idele group has
kernel consisting of $u$ for which $u_f=q1_f$ for some
$q\in\mathbb Q^\times$, with no condition on the nonzero real
component.  Modulo the rational subgroup this is the full embedded
$\mathbb R^\times$.  It includes the class of $(1_f,-1)$; in
$C_{\mathbb Q}\simeq\widehat{\mathbb Z}^{\times}\times
\mathbb R_+^\times$ that class has finite component $-1$.
Equivalently $(1_f,0)$ and $(-1_f,0)$ already represent the same point
of $Y_{\mathbb Q}$.  This is an explicit obstruction to the asserted
identification with $\widehat{\mathbb Z}^{\times}$.

Cyclotomic reciprocity sends $-1$ to complex conjugation, so quotienting
by it gives precisely the Galois group of the maximal real abelian
subextension.  A character descends to the quotient if and only if its
value at $-1$ is $1$.
\end{proof}

\begin{remark}
The proof of the unlabelled archimedean-fiber theorem in
\texttt{Jacobian.tex} itself first identifies the kernel as
$\mathbb R^\times$ and subsequently replaces it by
$\mathbb R_+^\times$.  The missing sign is therefore a concrete
quotient error, not an RH issue.  The archimedean stratum in
2401.08401, \texttt{sect3.tex}, item (2) following Proposition
\texttt{ergodic23}, already retains its quotient by $\pm1$.
An oriented version would require changing the construction and
checking its functoriality; it is not the stated $Y_{\mathbb Q}$.
\end{remark}

\begin{proposition}[AC-G14: prime idempotents and a non-idempotent class]
\label{AC-G14}
In the Jacobian monoid defined in \S\texttt{jacob}, each finite-prime
image $e_p=[\mathbb Z[1/p]]$ is a nonidentity idempotent.  Every monoid
homomorphism from this Jacobian to a group sends all $e_p$ to the
identity.  Thus no such homomorphism can retain the Abel--Jacobi
embedding of the distinct primes.

The whole Jacobian monoid is not an idempotent semilattice.  In the
nondegenerate-metric part, the rank-one group
\[
 L=\{x\in\mathbb Q:v_p(x)\geq-1\text{ for every prime }p\}
\]
has $[L\otimes_{\mathbb Z}L]\ne[L]$.
\end{proposition}
\begin{proof}
The localization multiplication map
$\mathbb Z[1/p]\otimes_{\mathbb Z}\mathbb Z[1/p]
\longrightarrow\mathbb Z[1/p]$ is an isomorphism.  This proves
idempotence.  It is not isomorphic to $\mathbb Z$: multiplication by
$p$ is surjective on $\mathbb Z[1/p]$ and not on $\mathbb Z$.
Likewise $\mathbb Z[1/p]$ and $\mathbb Z[1/q]$ for $p\ne q$ are
distinguished by $p$-divisibility.  If $f$ maps a monoid to a group and
$e^2=e$, cancellation in that group gives $f(e)=1$.  Equation
\texttt{extAJ} identifies these very idempotents as the prime images.

The displayed $L$ is an additive subgroup because
$v_p(x+y)\geq\min(v_p(x),v_p(y))$.  It is the directed union of the
cyclic groups $(1/n)\mathbb Z$ as $n$ ranges over squarefree positive
integers, ordered by divisibility.  Tensor products commute with
these directed unions, and multiplication identifies
\[
 L\otimes L\simeq
 \{x\in\mathbb Q:v_p(x)\geq-2\text{ for every prime }p\}=:L_2.
\]
Indeed a tensor of the two cyclic groups maps injectively to
$(1/(nm))\mathbb Z$, and the resulting directed system has injective
transition maps; every denominator with exponents at most two divides
such an $nm$.

Any additive isomorphism between two nonzero subgroups of
$\mathbb Q$ is multiplication by a nonzero rational number: tensor the
map with $\mathbb Q$ and use the one-dimensionality of both
rationalizations.  A fixed rational multiplier has nonzero valuation
at only finitely many primes.  At every other prime $p$ the greatest
allowed denominator exponent in its multiple of $L$ remains one,
whereas in $L_2$ it is two.  Hence no such multiplier identifies $L$
with $L_2$.
\end{proof}

\begin{remark}
This corrects the introductory description of the entire target as
``an idempotent semilattice'' in \texttt{Jacobian.tex}.  Its
idempotent boundary, including the extended Abel--Jacobi image, is a
semilattice.  The rank-one-group monoid defined in the body of the
paper contains further elements.  The same example corrects the
sentence in the proof of the generic-fiber theorem that every rank-one
subgroup of $\mathbb Q$ is isomorphic to $\mathbb Z$.  The claimed
generic fiber is nevertheless $C_{\mathbb Q}$: the finite part of an
\emph{idele} gives a cyclic fractional ideal, and conversely a cyclic
$L_a$ together with $a_\infty\ne0$ forces $a$ to be an idele.
\end{remark}

\subsection{Euler products and a quantitative consequence of A}

\begin{proposition}[AC-G15: the prime-circle Euler product of an abelian cover]
\label{AC-G15}
Let $L/\mathbb Q$ be a finite abelian extension with group $G$, and
take its cover $\gamma(L)$ in 2501.06560, Theorem \texttt{main}.
For each finite prime $p$ write $I_p\subset D_p\subset G$ for inertia
and decomposition, $e_p=|I_p|$, and
$f_p=|D_p/I_p|$.  The inverse image of the circle $C_p$ is a disjoint
union of $|G|/(e_pf_p)$ circles, each of flow length $f_p\log p$.
This includes the ramified primes.

Define a dynamical Euler product using exactly these finite-prime
circles, not every singular orbit of the full adelic monoid.  Then,
for $\Re s>1$,
\[
 Z_{\mathrm{circ}}(L,s)
 :=\prod_{p}\prod_{\mathcal C\subset\gamma(L)^{-1}(C_p)}
       (1-e^{-s\ell(\mathcal C)})^{-1}
 =\zeta_L(s)=\prod_{\eta\in\widehat G}L(s,\eta).
\]
Here $L(s,\eta)$ is the primitive Dirichlet/Artin function with its
correct inertia-invariant factors.  Under A every nontrivial zero of
this Euler product has real part $1/2$.  The construction of the
circles and the Euler-product identity are unconditional.
\end{proposition}
\begin{proof}
Lemma \texttt{Gacts} identifies the fiber over a point of $C_p$ as the
transitive $G$-set $G/I_p$, since the vanishing set is exactly $\{p\}$.
The universal mapping torus of Proposition \texttt{mappingtorus} is
\[
 p^{\mathbb Z}\backslash
 \left(\prod_{q\ne p}\mathbb Z_q^\times\times
 \mathbb R_+^\times\right).
\]
Passing to the finite associated cover gives the mapping torus of
translation by $\overline{\operatorname{Frob}}_p$ on $G/I_p$.
For ramified $p$ the element is defined in $G/I_p$, rather than in
$G$: the away-from-$p$ component of the diagonal number $p$ maps to
this class by cyclotomic reciprocity.  This also follows by embedding
$L$ in a cyclotomic field: its prime-to-$p$ roots of unity have
arithmetic Frobenius $\zeta\mapsto\zeta^p$, while the $p$-power
cyclotomic part contributes inertia.

Every cycle of the translation on $G/I_p$ has length $f_p$.  There are
$|G/I_p|/f_p=|G|/(e_pf_p)$ cycles.  Traversing the base once takes time
$\log p$, so the lifted circle has the asserted length.  These cycles
are also the cosets of $D_p$ in $G$, and therefore label the primes
$\mathfrak p$ of $L$ over $p$; their norms are $p^{f_p}$.
Multiplying $(1-p^{-f_ps})^{-1}$ over them proves the first equality
with the Dedekind Euler product, including bad primes.

For completeness let $t=p^{-s}$.  The regular representation of
$G/I_p$ decomposes as the characters $\eta$ of $G$ trivial on $I_p$.
The determinant of translation by
$\overline{\operatorname{Frob}}_p$ on this representation is
\[
 \prod_{\eta|_{I_p}=1}(1-\eta(\operatorname{Frob}_p)t)
 =(1-t^{f_p})^{|G|/(e_pf_p)}.
\]
This proves the character factorization locally, with no factor from
a character nontrivial on inertia.  Absolute convergence for
$\Re s>1$ justifies the products.  Analytic continuation supplies the
global identity.  A locates the nontrivial zeros of each primitive
factor; their finite product has no additional nontrivial zeros.
The sole pole at $s=1$ and the trivial zeros are not part of that
claim.
\end{proof}

\begin{remark}
The archimedean fiber is a fixed-point stratum rather than a finite
prime circle.  Gamma factors do not follow by adding it as an Euler
factor to this periodic-orbit product.  A trace formula including the
archimedean distribution, such as \texttt{weilexplicit}, is a
different, precisely specified identity.
\end{remark}

\begin{theorem}[AC-G16: square-root bounds for smoothed monodromy counts]
\label{AC-G16}
Assume A.  Fix the finite abelian extension $L/\mathbb Q$ of AC-G15,
and a positive integer $N$ divisible by every ramified prime and by
the conductors of all characters of $G$.  For
$\phi\in C_c^\infty((0,\infty),\mathbb C)$ and $g\in G$ define
\[
 \Theta_{g,\phi}(x)=
 \sum_{\substack{p\nmid N,\ k\geq1\\
                 \operatorname{Frob}_p^{k}=g}}
        (\log p)\phi(p^k/x),\qquad x\geq1.
\]
Then, as $x\to\infty$,
\begin{align*}
 \Theta_{g,\phi}(x)
   &=\frac{x}{|G|}\int_0^\infty\phi(t)\,dt
       +O_{L,N,\phi}(x^{1/2}),\\
 \sum_{\substack{p\nmid N\\\operatorname{Frob}_p=g}}
       (\log p)\phi(p/x)
   &=\frac{x}{|G|}\int_0^\infty\phi(t)\,dt
       +O_{L,N,\phi}(x^{1/2}\log(2x)).
\end{align*}
In particular, if $\phi\geq0$ is nonzero, every monodromy class occurs
in its dilated support for all sufficiently large $x$.  These are
fixed-extension estimates; no uniform conductor bound is asserted.
\end{theorem}
\begin{proof}
Put $\Phi(s)=\int_0^\infty\phi(t)t^{s-1}\,dt$.  If
$\operatorname{supp}\phi\subset[a,b]\subset(0,\infty)$, integration
by parts after $t=e^u$ shows that $\Phi$ is entire and, on every fixed
vertical strip, $\Phi(\sigma+iT)=O_A((1+|T|)^{-A})$ for every $A$.

Let $\eta$ be a primitive Dirichlet character arising from a character
of $G$; include the conductor-one trivial character.  For $c>1$,
absolute convergence of the Dirichlet series and Mellin inversion give
\[
 \Psi_{\eta,\phi}(x):=\sum_{n\geq1}\Lambda(n)\eta(n)\phi(n/x)
 =\frac1{2\pi i}\int_{c-i\infty}^{c+i\infty}
       -\frac{L'}{L}(s,\eta)\Phi(s)x^s\,ds.
\]
Move this contour to $\Re s=-1/4$.  The functional equation reduces
the logarithmic derivative on this line to one in $\Re s=5/4$ and
gamma logarithmic derivatives; it is
$O_\eta(\log(2+|\Im s|))$.  The Euler series is absolutely convergent
on the reflected line.  Consequently the new vertical integral is
$O_{\eta,\phi}(x^{-1/4})$.

Here are the convergence details for the contour passage.  The usual
zero-counting estimate for a fixed completed Dirichlet function is
$O_\eta(T\log(2+T))$, and the number in a unit-height band is
$O_\eta(\log(2+T))$.  Choose horizontal heights avoiding their
ordinates by a constant times $1/\log(2+T)$.  The logarithmic
derivative, from its local partial-fraction expansion and the gamma
term, is then bounded by a power of $\log(2+T)$ on the horizontal
segments.  The rapid decrease of $\Phi$ makes their integrals tend to
zero.  The same rapid decrease and the zero count imply absolute
convergence of the sum of residues.  These are the analytic facts
underlying the Dirichlet-sector specialization of equation
\texttt{weilexplicit} in \texttt{Jacobian.tex}.

There is a residue
$\delta_\eta x\Phi(1)$ from the pole at $1$, where
$\delta_\eta=1$ only for the trivial character.  Each nontrivial zero
$\rho$ contributes $-m_\rho\Phi(\rho)x^\rho$.  The only possible
trivial zero crossed is at $0$; its contribution
$-\operatorname{ord}_{s=0}L(s,\eta)\Phi(0)$ is a constant.  Thus
\[
 \Psi_{\eta,\phi}(x)
 =\delta_\eta x\Phi(1)
   -\sum_{\rho}m_\rho\Phi(\rho)x^\rho
   +O_{\eta,\phi}(1).
\]
A gives $|x^\rho|=x^{1/2}$.  Absolute convergence bounds the zero sum
by $x^{1/2}\sum_\rho m_\rho|\Phi(\rho)|$, a finite constant times
$x^{1/2}$.

Deleting the Euler terms at a fixed prime $p\mid N$ changes the sum
by $O_{p,\phi}(1)$, uniformly for large $x$: only integers $k$ in the
interval
$[\log(ax)/\log p,\log(bx)/\log p]$ can contribute, whose length is
the fixed number $\log(b/a)/\log p$.  This justifies using the
primitive functions while imposing the common unramified cutoff.

Character orthogonality in the finite abelian group gives
\[
 {\bf1}_{h=g}=\frac1{|G|}\sum_{\eta\in\widehat G}
                 \overline{\eta(g)}\eta(h).
\]
Apply it with $h=\operatorname{Frob}_p^k$.  The preceding estimate for
each primitive character proves the first formula, since
$\Phi(1)=\int\phi(t)\,dt$.

To remove repeated orbits, the contribution with $k\geq2$ is at most
\[
 \|\phi\|_\infty
 \sum_{p\leq\sqrt{bx}}(\log p)
       \frac{\log(bx)}{\log p}
 \leq \|\phi\|_\infty\sqrt{bx}\log(bx)
 =O_\phi(x^{1/2}\log(2x)).
\]
This elementary estimate suffices; no prime-number theorem is used at
this step.  It proves the second formula.  If $\phi$ is nonnegative
and nonzero its integral is positive, and the main term eventually
dominates the error, proving the final assertion.
\end{proof}

\begin{remark}
Under the time change $x=e^T$, the repeated-orbit error is
$O(e^{T/2})$ and the prime-only error is $O(Te^{T/2})$ in a fixed
length window of flow times.  This is a substantive A-consequence of
the existing geometric identification of monodromy.  It does not
produce a new nonabelian covering functor or GRH for Hecke characters
outside the Dirichlet sectors.
\end{remark}

\begin{proposition}[AC-G17: the trivial-cover exception]
\label{AC-G17}
Under Definition \texttt{deframi} of 2501.06560, the cover associated
to $G=\{1\}$ is unramified outside the empty set, and every
restriction above $C_p$ is trivial.  Hence Corollaries
\texttt{ramicor}, \texttt{monocor} and the archimedean assertion in
Theorem \texttt{main}(ii) require $G\ne\{1\}$.
For a nontrivial finite abelian $G$, their corrected conclusions hold:
the smallest ramification set in that definition is
$C(\chi)\cup\{\infty\}$, and infinitely many unramified prime circles
have nontrivial monodromy.  The obligatory $\infty$ in this particular
definition does not mean that the number field is nonreal.
\end{proposition}
\begin{proof}
For the trivial group the associated quotient is $X_{\mathbb Q}$
itself.  A singleton group acts freely and transitively on a singleton
fiber for every point; the implication in Definition
\texttt{deframi} therefore holds even with $S=\varnothing$.
There is no prime with $\chi(p)\ne1$.

Now suppose $G\ne1$.  Proposition \texttt{ramiprop}, whose proof uses
the stabilizer in Lemma \texttt{Gacts}, shows that a set containing
$\infty$ works exactly when it contains $C(\chi)$.  To see that
$\infty$ is necessary, let $S$ be a finite set not containing it and
put $T=S\cup C(\chi)$, a finite set of finite primes.  Take an adele
$a$ equal to $1$ at $T$ and $0$ elsewhere.  For every
$g\in G$, choose a residue class $r$ prime to the conductor with
$\chi(r)=g$.  The Chinese remainder theorem gives a positive integer
$n$ in that class which is prime to every additional prime of $T$.
Choose $w\in\widehat{\mathbb Z}^\times$ with $w_v=n$ at $v\in T$.
Then $aw=na$, so $w$ stabilizes $[a]\in Y_{\mathbb Q}$, and
$\chi(w)=g$ because $C(\chi)\subset T$.
Thus the stabilizer has image all of $G$, the fiber is a singleton,
and the nontrivial group does not act freely on it.  Yet the vanishing
set of $a$ is disjoint from $S$.  This proves failure of the defining
condition for every such $S$.

Choose any $g\ne1$ in $G$.  Surjectivity of $\chi$ supplies a reduced
residue class modulo its conductor mapping to $g$.  Dirichlet's
theorem gives infinitely many primes in that class, so Theorem
\texttt{mappingtorus1} gives nontrivial monodromy there.  This last
infinitude is unconditional; AC-G16 adds a quantitative estimate
under A.  The adele used to force $\infty$ also vanishes at many finite
places, so this ramification convention is not identical to the usual
real-place ramification of number fields.
\end{proof}

\subsection{Local perfectoid geometry and the tropicalization gate}

\begin{proposition}[AC-G18: a compatible tower is not its first coordinate]
\label{AC-G18}
Let $C$ be an algebraically closed perfectoid field of characteristic
zero and residue characteristic $p$.  In the identification of
\texttt{FF.tex}, Theorem \texttt{main}, the trivial local point is the
constant sequence $(1,1,\ldots)$, or the element $1$ of $C^\flat$.
It is not the whole locus $x_0=1$ in
\[
 C^\flat=\varprojlim_{x\mapsto x^p}C.
\]
That locus contains a nonconstant compatible cyclotomic tower, which
defines a nontrivial local point.  Consequently the prose instruction
``Removing the trivial point $x_0=1$'' after Theorem \texttt{main}
and in the matched-prime case of the universal-tilting subsection
must be read as removal of the all-ones sequence.
\end{proposition}
\begin{proof}
Choose primitive $p^n$-th roots $\zeta_{p^n}\in C$ with
$\zeta_{p^{n+1}}^p=\zeta_{p^n}$ and $\zeta_1=1$.
The sequence $x_n=\zeta_{p^n}$ is a point of the displayed inverse
limit and has $x_0=1$, but it is not the identity sequence.
Each $x_n$ is a principal unit: its reduction in characteristic $p$
satisfies $(\overline{x_n}-1)^{p^n}=0$, whence
$\overline{x_n}=1$.  Define
\[
 \chi(a/p^n)=\zeta_{p^n}^{a}\qquad(a\in\mathbb Z).
\]
Compatibility makes this a well-defined homomorphism on
$\mathbb Z[1/p]$.  Its kernel contains $\mathbb Z$, which is open in
the induced $p$-adic topology on $\mathbb Z[1/p]$.  It is therefore
continuous (indeed locally constant), and is nontrivial at $1/p^n$
for $n\geq1$ with a primitive root distinct from $1$.
Theorem \texttt{main} identifies it with a nonidentity element of
$1+\mathfrak m_{C^\flat}$.  Distinct compatible sequences remain
distinct in the inverse limit even when their zeroth coordinates
agree.  This proves the distinction.
\end{proof}

\begin{remark}
For a perfect characteristic-$p$ target, a compatible root sequence
is uniquely determined by its first coordinate, so the analogous
shorthand causes no problem there.  The issue arises when the
characteristic-$p$ formulation is transported to mixed
characteristic.  The quotient description of untilts in the stated
algebraically-closed setting uses
$(1+\mathfrak m_{C^\flat})\setminus\{1\}$, not a puncture defined by
the sharp/zeroth-coordinate map.  It also quotients by
$\mathbb Q_p^\times=\mathbb Z_p^\times p^{\mathbb Z}$; the compact
factor cannot be silently omitted.
\end{remark}

\begin{proposition}[AC-G19: the precise algebraic content of valuation profiles]
\label{AC-G19}
Suppose the Gauss valuations $v_s$ in \texttt{FF.tex},
\S\texttt{outlook}, have the asserted finite, concave, piecewise
affine profiles with integral slopes on the interval where they are
defined.  For $f\ne0$ set $\tau(f)(s)=-v_s(f)$ and put
$\tau(0)=-\infty$.  Then
\[
 \tau(fg)=\tau(f)+\tau(g),\qquad
 \tau(f+g)\leq\max\{\tau(f),\tau(g)\},
\]
pointwise.  Equality in the second formula holds at every parameter
where $v_s(f)\ne v_s(g)$.
The map $\tau$ from the ring itself is not a semiring homomorphism to
the max-plus functions.  Consequently the profile assertion alone
does not define the claimed functor of ringed or semiringed spaces;
the additional construction and its gluing and descent remain a
mathematical gate.
\end{proposition}
\begin{proof}
The first formula is the multiplicativity axiom of a valuation with
its sign reversed.  The ultrametric inequality
$v_s(f+g)\geq\min(v_s(f),v_s(g))$ gives the second formula.  If, say,
$v_s(f)<v_s(g)$ and the inequality were strict, applying it to
$f=(f+g)-g$ would give
$v_s(f)\geq\min(v_s(f+g),v_s(g))>v_s(f)$, a contradiction.

Take any nonzero $f$ and $g=-f$.  Since $v_s(-1)=0$, the two profiles
agree and are finite.  Their max is that finite profile, whereas
$\tau(f+g)=\tau(0)=-\infty$.  Thus preservation of addition fails.
This counterexample works even with $f=1$, and is independent of RH.
Negating the asserted concave profiles does give convex profiles with
integral slopes on the original parameter interval.  It does not
repair the failure of additive functoriality.
\end{proof}

\begin{remark}
Two additional scope checks are necessary before interpreting these
profiles as global sections of the scaling-site sheaf.  First the
valuation parameter in the source proposition ranges over
$(0,\infty)$, whereas the scaling site's underlying half-line
includes $0$; an extension at $0$ must be established for the chosen
class of functions.  Second a profile on the half-line is not by
itself a section descending to a periodic quotient $C_p$; the
appropriate Frobenius transformation law and descent must be checked.
The source's final paragraph explicitly leaves the eigenspace
comparison and descent question open.  AC-G19 refutes only the direct
semiring-homomorphism inference; it does not rule out a separately
constructed tropicalization via valuations, ideals, or other
appropriate categorical data.
\end{remark}

\begin{proposition}[AC-G20: exact limitation of the abelian monodromy data]
\label{AC-G20}
The mapping-torus actions in 2401.08401 and 2501.06560 factor through
abelianized etale fundamental groups.  Every finite quotient of those
specified actions is abelian.  They cannot, without changing the
construction, provide the full nonabelian Galois action of an
$S_3$-extension of $\mathbb Q$.

Likewise the group $\mathbb Q_p^\times$ acting by continuous additive
automorphisms of the completed stalk in \texttt{FF.tex} is an
abelian local symmetry group.  Its notation as a local Weil group
does not identify it with the full Weil group of $\mathbb Q_p$.
\end{proposition}
\begin{proof}
The fundamental groups in the cited mapping-torus theorems are
explicitly abelianized, and their displayed realizations are products
of $\mathbb Z_q^\times$.  Their subgroups and their group quotients
are abelian.  The image of a homomorphism from an abelian group is
abelian, whereas $S_3$ is not: for example $(12)(23)\ne(23)(12)$.
Thus no surjective action with that Galois group can factor through
the stated source group.  This argument concerns the specified
class-field construction, not every possible automorphism group of
the underlying topological space.

For the local assertion, a continuous additive endomorphism
$f:\mathbb Q_p\to\mathbb Q_p$ satisfies $f(a)=a f(1)$ first for
$a\in\mathbb Q$ by additivity and division, and then for all
$a\in\mathbb Q_p$ by density and continuity.  It is invertible
exactly when $f(1)\ne0$.  Hence
$\operatorname{Aut}_{\mathrm{top,add}}(\mathbb Q_p)
=\mathbb Q_p^\times$, an abelian group.  Local reciprocity concerns
the abelianization of the Weil group; a nonabelian Weil representation
does not become a representation of this stalk-symmetry group merely
by naming it $W_p$.
Indeed the full local Weil group already has nonabelian tame finite
images: choose an integer $m>1$ prime to $p$ with
$p\not\equiv1\pmod m$, and let $f$ be the order of $p$ modulo $m$.
The unramified extension of degree $f$ contains $\mu_m$; adjoining an
$m$-th root of a uniformizer gives a finite Galois extension with
inertia $C_m$ and Frobenius acting on it by $\tau\mapsto\tau^p$.
This action is nontrivial, so the image is nonabelian.  Such an $m$
exists by choosing a prime divisor of neither $p$ nor $p-1$.
\end{proof}

\begin{proposition}[AC-G21: the finite character product uses a pro-ring, not the asserted Ind-ring]
\label{AC-G21}
In \texttt{Jacobian.tex}, \S\texttt{sectgeompic}, the remark
``The hidden coproduct structure'' proposes to regard
$\mathbb Q/\mathbb Z$ as the direct limit of the rings
$R_n=\mathbb Z/n\mathbb Z$, identifying $1/n$ with $1_n$,
and to take the direct limit of the maps
$\Delta_n:R_n\simeq R_n\otimes_{\mathbb Z}R_n$.
With those additive-group identifications this is not a diagram of
rings, even if nonunital ring maps are allowed, so that proposed
Ind-ring coproduct is not defined.

The preceding levelwise character-product construction in equation
\texttt{tensor} remains valid: its compatibility uses the inverse
system of unital ring reductions $R_N\to R_n$ for $n\mid N$.
It can be proved without a coproduct on $\mathbb Q/\mathbb Z$.
\end{proposition}
\begin{proof}
The inclusion of cyclic subgroups
$\langle1/n\rangle\subset\langle1/N\rangle$ for $N=kn$
corresponds, under the specified identifications, to the additive map
\[
 j_{n,N}:R_n\longrightarrow R_N,\qquad
 j_{n,N}(a)=ka\pmod N.
\]
For $n=2$, $N=4$, this sends $1$ to $2$, but
$j_{2,4}(1)^2=0$ and $j_{2,4}(1^2)=2$ in $R_4$.
Thus it is not multiplicative.  Nor can an alternative unital ring
map $R_n\to R_N$ exist when $n\mid N$ and $N>n$: applying such a
map to $n1=0$ would give $n1_N=0$, contrary to $N\nmid n$.
Consequently the direct limit in the remark is not an Ind-object of
rings with the displayed transition maps.

The additive obstruction is stronger.  If $a\in\mathbb Q/\mathbb Z$
has order dividing $n$, write $b=nb'$ using divisibility.  Then
$a\otimes b=na\otimes b'=0$.  Every pure tensor vanishes, so
$(\mathbb Q/\mathbb Z)\otimes_{\mathbb Z}
(\mathbb Q/\mathbb Z)=0$.  Hence every bilinear multiplication on
that additive group is zero; it cannot restrict to the stated
multiplications on its finite cyclic subgroups.

Here is the corrected compatibility argument for the actual character
product.  Let $\rho_i$ be roots for $L_i$, and encode the value
$\rho_i(1/n)$ by an additive map
$u_{i,n}:L_i\to R_n$: its value as a $\mathbb Q/\mathbb Z$-valued
character is $u_{i,n}(x)/n$.  The relation
$\rho_i(1/n)=k\rho_i(1/N)$ says precisely
$u_{i,N}\bmod n=u_{i,n}$.
Define the additive map $w_n:L_1\otimes L_2\to R_n$ by
\[
 w_n(x\otimes y)=u_{1,n}(x)u_{2,n}(y).
\]
It exists by the universal property of the tensor product, since the
displayed pairing is bilinear.  The reduction $R_N\to R_n$ is a
ring homomorphism, and therefore $w_N\bmod n=w_n$.  Consequently
\[
 \rho(k/n)(z)=\frac{k\,w_n(z)}{n}\pmod{\mathbb Z}
\]
is independent of the chosen representation of $k/n$ and defines a
group homomorphism
$\mathbb Q/\mathbb Z\to
\operatorname{Hom}(L_1\otimes L_2,\mathbb Q/\mathbb Z)$.
For example replacing $k/n$ by $k(N/n)/N$ changes the value to
$k w_N(z)/n$, which agrees modulo $\mathbb Z$ with the displayed
one; common denominators then prove additivity.

If the roots come from frames
$\xi_i:L_i\to\widehat{\mathbb Z}$, their residue maps are exactly
$u_{i,n}$.  The product frame
$x\otimes y\mapsto\xi_1(x)\xi_2(y)$ reduces to $w_n$ at every
level, so its dual has precisely this root.  This proves the
compatibility needed in Lemma \texttt{dualtensor} and preserves the
frame-product argument.  The projections form a pro-ring; no
direct-limit ring structure or ordinary tensor coproduct on
$\mathbb Q/\mathbb Z$ is required.
\end{proof}

\subsection{What remains a genuine geometric gate}

The new papers supply substantial existing geometry: a Picard monoid,
its arithmetic Abel--Jacobi map, sheaves and local points, finite
abelian coverings, mapping tori, and local perfectoid and complex
quotients.  Under A the Dirichlet spectral sectors acquire their
critical-line location; AC-G15 and AC-G16 give precise Euler-product
and prime-orbit consequences.  None of these implications alone
constructs a geometric square with a surface Riemann--Roch theorem,
a cohomology whose polarized trace is the global explicit formula, or
a motive uniquely determined by the analytic Hodge inequality.
The trace regularization, comparison of the local constructions with
one global cohomology, Frobenius-compatible tropical descent, and
extension beyond abelianized monodromy are separate tasks with
separate hypotheses.  The concrete corrections AC-G13, AC-G14,
AC-G17, AC-G18, AC-G21 and the categorical obstruction AC-G19 must be resolved
before those supplied source formulations can be used as premises.

The Grothendieck crosswalk remains the exact one recorded in the
original module: Montpellier catalogue cotes 2, 10, 11, 14, 16 and 19
provide themes (Riemann--Roch, tensors, correspondences,
Hodge--Tate--Lefschetz, and topoi), not verified scan-level premises
for any theorem here.  No additional archival statement or scan
number is attributed in this follow-up.  The use of localization,
tensor products, inertia and quotient sheaves is explicit in the
proofs, and does not supply an unverified archival quotation.
